\section{BitLocker}

Windows-Festplattenverschlüsselung.
Zum Entschlüsseln wird der Volume Master Key (VMK) benötigt.

\subsection{VMK aus SPI-Dump extrahieren}

Hardware-Angriff über Logic-Analyzer-Daten:

\begin{lstlisting}[language=bash]
# YAFFS in Firmware suchen
binwalk firmware.bin | grep -i yaffs

# VMK mit ARNE-Tool extrahieren
./arne -i digital.csv \
  -k "1=SCLK,2=CS,3=MOSI,4=MISO" -o key.vmk
xxd key.vmk
\end{lstlisting}

\subsection{BitLocker-Volume entschlüsseln}

\begin{lstlisting}[language=bash]
mkdir -p /tmp/bitlocker /tmp/windows

# Mit VMK-Key entschluesseln
dislocker /dev/loopXpN -K key.vmk -- /tmp/bitlocker

# NTFS dirty flag beheben
ntfsfix -d /tmp/bitlocker/dislocker-file

# Entschluesselte Partition loopen
losetup --partscan --find --show \
  /tmp/bitlocker/dislocker-file
\end{lstlisting}

\section{dmsetup — Partitionen zusammenführen}

Nach der Entschlüsselung müssen die Partitionen zu einem virtuellen Gesamtdisk zusammengeführt werden:

\begin{lstlisting}[language=bash]
# dmsetup.txt (Loop-Nummern anpassen):
# 0 239616 linear /dev/loop37 0
# 239616 248416256 linear /dev/loop38 0
# 248655872 1411072 linear /dev/loop37 248655872

cat dmsetup.txt | dmsetup create merged
\end{lstlisting}

\section{QEMU-Virtualisierung}

Entschlüsseltes Image als VM starten:

\begin{lstlisting}[language=bash]
qemu-system-x86_64 -bios /usr/share/ovmf/OVMF.fd \
  -cpu Skylake-Client-v3 -enable-kvm -usb \
  -hda /dev/mapper/merged -smp 4 -m 4096
\end{lstlisting}

\section{VeraCrypt / TrueCrypt}

Container erkennen durch hohe Entropie:

\begin{lstlisting}[language=bash]
binwalk -E verdaechtige_datei
# Hohe Entropie (>7.5) = wahrscheinlich verschluesselt
\end{lstlisting}

Container öffnen (falls Passwort bekannt):

\begin{lstlisting}[language=bash]
veracrypt --text --mount container.vc /mnt/container \
  --password="PASSWORT" --non-interactive
\end{lstlisting}

\section{Aufräumen (LIFO)}

\begin{lstlisting}[language=bash]
dmsetup remove merged
losetup -d /dev/loop38
umount /tmp/bitlocker
losetup -d /dev/loop37
umount /tmp/ewf
\end{lstlisting}
